\chapterstart{综合讨论}
\label{sec:discussion}

\subsection{有限经验的价值与上下文依赖}

本研究的核心认识是：有限经验的学习价值，取决于它在预测中提供了什么信息关系，以及训练是否要求模型利用这些关系。文本内容、相对位置、表达形式与预测目标共同构成一次学习机会。同样材料分窗后来源效应减弱，说明关系能否被同时访问具有实际作用；重复与重述对不同目标产生不同效果，则说明训练关系会塑造后续的信息使用方式。

这一认识为数据处理提供了更具体的判断依据。扩展覆盖、精简表达和重复练习分别改变内容范围、预测线索和学习频次；选择方法时，可以检查目标能力依赖什么信息、这些信息是否进入预测窗口、监督是否实际训练这种使用。本研究进一步在既有文本上改变遮盖与监督，在相同累计曝光下取得优于普通续训的结果，将关系研究转化为可采用的训练设计。

相同累计曝光下的方法改善，是本文关于学习效率的直接实践证据。与拟合数据、计算资源和性能关系的标度研究\citep{muennighoff2025scaling}相联系，本研究补充了机制与设计层的解释：怎样组织有限经验，怎样检验学到的计算，以及怎样在后续训练中使用这些认识。跨数据量和模型规模的效率变化，可据此进一步建立比较曲线。

\subsection{新能力与已有功能的共存}

新能力的形成需要改变模型，已有功能的保存则需要约束其中一部分变化。二者应按使用条件协调：在新的关系预测任务上允许学习发生，在普通语言输入上减小不必要的预测偏移。实验中，密集遮盖输入上的强保持压低了新学习收益；普通输入保持则在保留大部分新预测改善的同时，减小了普通输入上的损失变化。

能力保持还包含复用范围的问题。受控任务中，熟悉输入的准确率与未见符号上的关系处理出现明显分离；调整监督分配后，未见符号的表现改善，内部信号干预也改变了回答。结合已有结构性上下文学习研究\citep{anand2025dual}，这些结果说明，保持评价需要同时观察熟悉行为和新输入对已学计算的使用。

普通输入保持将这些认识用于真实模型，并在两个续训种子上进一步改善完整评测。该结果验证了完整方案的实用性；保持输入、有效约束强度、更新幅度和额外呈现的各自贡献，是下一步可直接操纵和比较的变量。

\subsection{Qiushi Engine 的研究组织与积累}
\label{sec:research_organization}

Qiushi Engine 是团队自主研发的分层多智能体科研系统，其基本架构与前期真实光学平台研究见 Yang 等\citep{yang2026optical}。核心研究智能体分工承担研究规划、方法构造、实验执行和批判性审阅，支撑子智能体提供资料检索、辅助探索与独立核查。知识系统组织文献、数据和方法资料，记忆系统保存研究进展、实验依据及已形成的方法经验，使后续研究能够继续使用这些积累。各部分共享科研计算基础：团队开发的执行管理与工具接口连接算力和存储等基础设施（Infrastructure）、科学计算环境（Environment）及程序运行支持（Runtime）；AI Lab 为模型生成、训练、评测和机制实验提供可调用的科研工具。

Qiushi Engine 主导了从问题形成到成果整合的连续研究，具体工作随科学问题而展开：查阅文献、构造数据、提出模型与学习方法、实现实验、比较结果、分析机制，并据此决定下一项研究。紧凑重述的早晚排序不同，引出预算与学习阶段的比较；连贯与打散输入的差异，引出关系组织研究；熟悉与未见符号的表现分离，引出监督和保持实验；固定目标后仍观察到遮盖作用，又修正了对监督数量的早期解释。

研究的连续性体现在方法技术与科学认识的共同发展。紧凑重述既改变了训练数据，也为比较不同表达关系提供了可控制的材料；残差增量结构既用于改进模型，也让后续实验能够固定基座、比较不同的学习目标。已有的训练和评测实现被继续使用，关于方法为什么有效的解释则随新实验不断修正。第三阶段因此能够在第一代模型与数据的基础上，进一步设计输入、监督与保持的配合方式。

\begin{table}[htbp]
\caption{研究结果与后续科学决策的联系。观察到的现象形成新问题，并推动相应的对照、机制分析和训练设计。}
\label{tab:research_decisions}
\small
\begin{tabularx}{\linewidth}{@{}Y Y Y@{}}
\toprule 改变判断的观察 & 提出的新问题 & 后续实际研究\\\midrule
紧凑视图的早期劣势与晚期优势 & 数据价值是否依赖预算和能力状态 & 固定预算替换、晚期轨迹与跨结构比较\\
连贯片段与打散材料产生差异 & 相同词语究竟练习了哪种关系 & 重复、重述、错源及分窗对照\\
熟悉行为恢复而未见调用弱 & 能力恢复的是表现还是调用范围 & 监督分配、信号抹除与跨位置搬移\\
密集输入固定稀疏目标仍有效 & 收益是否必须增加监督数量 & 输入遮盖与目标集合分别比较\\
普通与密集保持强度不同 & 保持配置如何影响获得与保存 & 共同目标梯度比较及完整方案评测\\
\bottomrule
\end{tabularx}
\end{table}

这些联系体现了长期自主研究的具体价值：已有结果持续改变下一项研究决策，已有材料持续支持新的实验。Qiushi Engine 将方法构造、实验实现、结果分析和解释修正连接起来，使一次模型突破成为后续科学认识与方法创造的起点。

自主科研中的想法生成、实验迭代、分析与写作已有系统研究\citep{lu2026automation}。本项目提供了有限数据语言学习中的三阶段案例：从前沿结果中提出科学问题，再让得到的认识改变后续模型训练。开放这条研究链及其材料，可以支持社区研究长期知识积累如何影响自主科研的质量与效率。

\subsection{Research RSI：科研过程的递归自我改进}
\label{sec:rsi}

长程自主研究的积累，不止于保存已经找到的最佳模型。对数据价值的判断、对学习机制的认识、可组合的方法技术和能够区分解释的实验，也会成为下一阶段的研究基础。Qiushi Engine 的三个阶段展示了这些成果怎样共同发展：实践扩展可行的方法空间，机制研究辨明方法的作用条件，新的模型创造则检验这些认识能否转化为实际改进。

我们用 \textbf{Research RSI}（Recursive Self-Improvement of the Research Process，科研过程的递归自我改进）描述其中的递归关系。持续开展研究的 Qiushi Engine 将自身形成的科学认识、方法创新与实验经验重新用于后续研究，更新问题选择、实验策略、方法构造和证据判断；新的实践再检验并修正这些积累。这里直接更新的是系统开展科研所依据的知识与方法，而非研究代理的参数或代码。\textbf{一次研究不仅改变系统得到了什么，也改变系统下一次怎样做研究。}

本项目中的改变同时涉及研究认识、实验工具和训练方法。紧凑重述的早晚期差异，使数据选择从静态排序转向预算、替代对象和学习阶段的共同比较；残差增量结构从模型改进方法发展为固定共同基座、研究功能变化的实验工具；关系与能力复用实验进一步促成输入、监督和保持的分别设计。它们没有被压缩成一份最优参数清单，而是以带有作用条件、对照结果和修正记录的方法继续参与研究。表\ref{tab:research_decisions}与表\ref{tab:bridge}分别列出判断的演进及其进入训练的具体方式。

新模型取得收益之后，研究仍继续检验设计依据。固定目标后的遮盖比较，修正了“增加监督数量带来改善”的早期解释；共同位置上的梯度比较，则表明相同保持系数可以产生很不同的实际约束。因而，下一轮研究既可使用有效方案，也能利用更准确的比较条件。科研笔记、实验规划、程序与结果保存了这一过程，使读者能够看清一项判断从何而来、怎样进入方法、又被哪些证据改变。这是本案例可解释性的主要来源。

相关研究已经从不同对象探索自我改进：STOP 将程序改进器用于自身\citep{zelikman2024stop}，Darwin G\"odel Machine 修改智能体代码并以任务评测筛选改进\citep{zhang2026dgm}，Reflexion 将语言反馈保存在记忆中供后续尝试使用\citep{shinn2023reflexion}。Co-Scientist 通过假设生成、辩论、演化和持久记忆组织自我改进循环，并与科学家合作完成生物医学实验验证\citep{gottweis2026coscientist}；近期综述也将科研过程本身列为一种改进对象\citep{chen2026recursive}。Qiushi 的具体贡献在于，把真实前沿模型的创造、学习原理的发现、原理指导的后续模型设计及完整比较连接在同一研究计划中，并保留各阶段知识与方法如何发展、结合和修正的实证材料。

本案例支持同一 BabyLM 项目内的递归研究闭环：研究积累改变了后续设计，真实模型对照验证了其中具体方案的收益，进一步实验又修正了认识。它尚不等同于跨任务的一般科研能力已持续提升；这一更广问题需要在独立研究目标上比较知识复用对实验选择、研究成本与成果质量的影响。三阶段的意义，是使前沿突破既成为一项研究成果，也成为后续研究可以继承、检验和发展的知识来源。

\subsection{可复用成果与后续检验}

本研究留下了三个相互联系的贡献：关系类型与预测窗口影响上下文使用的实验发现，监督分配与内部干预揭示能力复用条件的机制研究，以及在真实模型上得到重复改善的训练方案。经验压缩、关系锚点、身份匹配和测量控制等独立研究则扩展了可供继续检验的方法范围。

这些成果在不同层次上提供了可继续研究的起点。自然文本对照揭示关系组织的选择性作用，受控任务通过内部干预解释新输入怎样使用已学计算，真实训练则验证输入、监督与保持设计的实际收益。三类证据相互补充，共同形成从学习现象到方法创造的联系；它们各自的对照范围已在相应实验中说明。独立母模型、不同预算和文档隔离条件下的检验，可以进一步研究这种联系的适用范围。

分项分析进一步说明了方法的实际价值。相对母模型，GlobalPIQA 贡献两个完整方案 Overall 增量的约 74\% 与 79\%；相对普通续训，最大的正向贡献则来自实体跟踪。加入保持方案后的附加收益来自其他分项的净变化，GlobalPIQA 在这一步保持不变。固定模型后的两种微调种子比较还表明，完整方案的 (Super)GLUE 优势不只见于一次微调。区分这些参照，才能解释每一步训练设计改善了什么，而不把一种比较的收益结构套用于所有结论。

当前完整结果来自同一母模型的两次续训，完整方案增加了保持性文本呈现。GlobalPIQA 仅包含 203 题，两个续训模型使用相同评测集合，逐题配对的不确定性区间尚未随本文提供。因而，本研究报告的是已有对照下的综合净改善及具体分项变化；训练随机性、评测样本不确定性与跨底座迁移需要分别衡量。

反例与修正也是可复用知识。任务捷径检查、共同参数对齐和文本重合分析，分别排除了答案构造、初始化和样本选择中的混淆。它们帮助后来者建立更可靠的对照，并明确哪些早期解释已被后续结果替代。

成果的使用因此包括模型和研究方法两个层次。两代模型提供实践起点；数据构造、目标定义、对照设计、分项结果和研究笔记说明方法如何形成、怎样继续检验。报告与配套仓库共同组织这些材料，具体文件入口和当前可用范围见附录\ref{app:assets}。
